\section{Discussion}
\label{sec:discussion}

\subsection{AE4 supplies chloride; secretory pathways expose its availability}

The principal result is a conditional mechanism with an explicit mass balance. Wild type begins with more intracellular chloride and retains AE4 mediated chloride loading. When the two genotypes receive matched non-AE4 supply and share apical conductance parameters, knockout exports less chloride and secretes less fluid. A substantial, persistent deficit therefore does not require an AE4 specific apical multiplier in this construction. The reservoir identity accounts for the additional export without invoking an unmeasured source of chloride.

This is stronger than a verbal claim that AE4 maintains a chloride store, but narrower than a complete causal identification of the native knockout phenotype. The initial state is projected using measured chloride and pH, not obtained as a chronic genotype equilibrium. Matched supply is prescribed, not derived from transporter regulation. Projected TIC and potassium differ between genotypes, so the experiment is not an isolated manipulation of chloride alone. The result establishes sufficiency of a particular chloride availability construction within the conserved network, while leaving its biological parameterisation open.

The distinction between a reservoir and a finite store is also important. The model does not empty the entire initial wild type advantage. That difference is larger at 600 s than at onset, and AE4 has supplied additional chloride throughout the interval. The correct explanation is continuing replenishment together with differential availability, not an initial packet of chloride that is simply spent until exhaustion.

\subsection{What changed since the original model}

Our 2018 article provided a plausible network explanation for the striking AE4/AE2 distinction \citep{VeraSiguenza2018}. Its cation dependent AE4 cycle altered intracellular sodium, potassium, NHE1 and pump activity, and predicted weak NKCC activation after exchanger deletion. It should not retrospectively be described as the strongly compensating model: that behaviour emerged in later reconstructions with different transport laws and state responses.

The present analysis changes the evidential status of the explanation rather than negating the original biological observation. The original concentration formulation contained acid--base chemistry, but a successful secretion trace alone did not identify whether sustained alkalinity supply and the partition of transporter fluxes were biologically correct. The amount formulation makes those requirements easier to test. It also exposes that replacing a constitutive law can change the apparent role of AE4 even when the conservation equations and transported stoichiometry remain correct.

The final comparison addresses a further difference. Acute removal at a common wild type state and secretion from established knockout cells are not the same perturbation. A common onset deliberately removes the measured genotype chloride difference. Conversely, imposing that difference does not explain how chronic knockout cells acquired or maintained it. Neither comparison should silently be substituted for the other. Their contrast identifies the need for a future independently constrained account of chronic state formation, not a licence to tune resting conditions against the secretion endpoint.

\subsection{The pitfalls that changed the mechanistic interpretation}
\label{sec:pitfalls}

\paragraph{Conservation is necessary, but not a sufficient biological test.}
An electroneutral source can still drive pH outside the physiological domain. Carbon dioxide transport supplies carbon without net alkalinity, bicarbonate contributes one equivalent of each, and carbonate contributes one carbon but two alkalinity equivalents. Conflating these quantities can make an apparent chloride supply look sustainable when its acid--base budget is not. In the frozen reconstruction, an NBC-like input supplies the additional alkalinity needed for substantial productive AE4 flux. The appropriate inference is conditional on that architecture and its allowed fluxes, not universal necessity of a named transporter. Similarly, a passing 600 s numerical check does not rule out slow modes or later failure.

\paragraph{A plausible compensating law can conceal the phenotype.}
The Palk/Benjamin reconstruction allows changes in intracellular substrates after AE4 removal to raise NKCC chloride uptake substantially. Within that model the compensation is real; the mistake would be to treat the mathematical response as experimentally established merely because the rate law is familiar. The isolated NKCC experiment found no detectable genotype difference under a specific depletion and readdition protocol \citep{Pena2015}. Its conditions are not normal whole cell stimulation, and the reduced Palk law does not reproduce the external substrate manipulation faithfully. Thus neither the large simulated compensation nor exact equality of physiological fluxes is proved by that assay. The final matched-supply comparison is useful precisely because it makes the alternative assumption explicit.

\paragraph{An algebraic cancellation is not zero biological sensitivity.}
Equation~\eqref{eq:equal-routing-identity} removes the explicit AE4 term for equal routing, but AE4 can change pump and NHE1 activity through the state. The identity does not make secretion independent of AE4, nor does it identify a globally optimal intervention or a universal minimum coupling strength. Its value is diagnostic: it explains why searching source coefficients alone may fail to control sustained export. The chloride identity in \cref{eq:reservoir-general} avoids this routing restriction and remains useful for finite time comparisons.

\paragraph{Source stoichiometry and a kinetic law are different commitments.}
A proposed cycle can pass charge accounting while an inherited scalar law drives it against its own affinity. The seven-class panel therefore tests a specified source replacement, not every possible kinetic realisation of the molecular proposal. Its physiological failures and small or reversed deficits are informative exclusions within that construction. They cannot be promoted to a refutation of the biochemical measurements. Full consistency would require class specific kinetics with compatible reversal, substrate dependence and capacity, constrained independently of secretion.

\paragraph{Opening a volume sensitive pathway requires an initiating perturbation.}
The earlier swelling-only VRAC construction attached a positive swelling gate to an AE4 knockout core with no beta driven upstream change. At an exact unchanged resting state, such a gate cannot initiate its own activating perturbation. This excludes that circular initiation mechanism, not every role of volume sensitive channels. Adding beta driven NKCC loading changes the premise and permits local chloride loading and swelling. However, local initiation does not identify the conductance needed for sustained secretion.

The subsequent independently targeted swelling calculation illustrates the distinction. At the two frozen conductance endpoints, mean swelling was approximately 5.75\% and 2.44\%, rather than the prescribed 12.5\% target. The residuals did not bracket a root, so no accepted conductance or quantitative beta-protocol prediction followed. Two endpoints do not prove global absence of a solution or monotonicity. The scientifically useful result is the unresolved identification, retained without extra searches or phenotype based rescue.

\paragraph{Calibrated sufficiency must not become an inferred molecular mechanism.}
Equation~\eqref{eq:calibrated-coupling} demonstrates that an added interaction can generate a large deficit, but its coefficient derives from the phenotype. Agreement is therefore not an independent validation. The final reservoir result shows that such a multiplier is not required for a substantial deficit in the stated alternative construction. Conversely, it does not exclude additional beta conditioned interactions in the gland. The two models are not additive components, and their percentage effects cannot be summed or subtracted to estimate an unobserved pathway.

\subsection{Catalán 2015: a shared demand pathway, not necessarily an AE4 specific channel}

The 2015 TMEM16A deletion experiment separates muscarinic calcium dependent secretion from an IPR responsive route \citep{Catalan2015}. Persistence of the latter, its swelling association and sensitivity to DCPIB/NPPB support considering an auxiliary anion pathway. They do not establish a direct AE4-to-channel interaction or a unique molecular channel identity. Our shared auxiliary conductance is an effective representation of demand, not a claim that the entire measured whole cell current is apical or that its conductance follows directly from a holding voltage.

The final model permits a different location for the genotype dependence: substrate availability rather than channel abundance. Identical conductance parameters act on different chloride concentrations and membrane potentials, so their currents differ. The modest effect of varying auxiliary strength on the relative deficit indicates that the shared pathway is not the adjustable source of the phenotype in this comparison. Indeed, a similar deficit persists when the added conductance is zero. This supports chloride availability as a general secretory constraint but does not explain the experimental selectivity of beta containing uptake responses.

\subsection{Catalán 2025: molecular constraints do not uniquely fix a gland model}

The human AE4 mutagenesis and molecular dynamics results identify a putative region coordinating cations and bicarbonate around the TM3--TM10 interface \citep{Catalan2025}. The different sodium and potassium responses of the T448I/T756A double mutant are especially important: cation participation is not a dispensable label on an otherwise specified chloride exchanger. Nevertheless, coordination and mutant activity do not uniquely determine whole-cycle stoichiometry, the sodium/potassium flux split at native concentrations, or the kinetic response during salivary stimulation.

The original concentration weighted allocation, the later equal allocation and the final retained donor allocation are therefore effective model choices, not rival direct measurements of the protein. Our source-class panel makes part of that uncertainty explicit, while also demonstrating the danger of combining a proposed source vector with an unchanged kinetic law. The reservoir identity does not resolve the molecular cycle; it provides a conservation constraint that any admissible cycle must satisfy once its actual chloride input has been specified. This is the useful connection between the molecular study and the present article: molecular evidence constrains possible transport, whereas the full cell budget determines how that transport can influence secretion.

\subsection{Why approximately 35\% is context rather than a mathematical constant}

The experimental deficit is an observation with sampling uncertainty, protocol dependence and a different spatial scale from a single model cell. The reported standard error is not a confidence bound and cannot define an acceptance threshold by subtraction from the mean. A 20\% conditional prediction is not automatically a failure because it differs from 35\%, just as a fitted 35\% result is not automatically validated.

The relevant achievement is a substantial reduction of the correct direction, persisting across the declared observation windows without selection of an AE4 specific channel coefficient against that outcome. The relevant remaining weaknesses are mechanistic: imposed onset states, matched supply, unmeasured projected quantities and incomplete stimulus specificity. It would be equally unjustified to claim that the model explains $20/35$ of the biological mechanism or that the numerical difference identifies a missing 15 percentage point pathway. Different constructions, observation maps and baselines do not support that arithmetic as causal attribution.

\subsection{Limitations and discriminating observations}

The source chloride and pH measurements constrain onset, but do not identify all intracellular carbon species, potassium, volume or their covariance. The central projection substantially changes TIC between genotypes. The reported alternate projection is a useful sensitivity, not proof that all admissible onsets agree. Likewise, the reported extra-supply cases suggest attenuation without sign reversal, but cannot replace measurements of non-AE4 fluxes under matched native stimulation. Their original arrays require deposition before those sensitivity claims are fully reproducible from this revision's source package.

The reported CCh only deficit and IPR pH failure remain substantive qualifications. A CCh only fluid deficit cannot be compared directly with preserved SPQ uptake: a fluorescence recovery slope, a concentration derivative and gland fluid output have different observation maps. For example, $\dd[\Cl]_i/\dd t=(\dot n_{\rm Cl,i}-[\Cl]_i\dot V_i)/V_i$ includes dilution. Nevertheless, retention of much of the model deficit without beta stimulation means protocol specificity has not been established. A trajectory that crosses the pH limit near 488 s has no accepted 600 s endpoint, regardless of what its unvalidated continuation might suggest.

No final reservoir AE2 knockout or chronic 5\% AE4 comparison was performed. The original AE2 specificity result and exact nesting properties of the calibrated benchmark must therefore not be transferred to the final construction. The single cell approximation, effective signalling inputs and absent ductal model also limit quantitative interpretation. The strict conservation identity does not remove these biological limitations.

The most discriminating additional observations would separate supply from demand under the same stimulation protocol: genotype specific chloride amounts and cell volume, non-AE4 chloride uptake, pH and bicarbonate handling, and apical current with an identified driving force. Such observations could determine whether compensation is genuinely constrained and whether the measured reservoir difference survives without imposing it. This is a narrower and more informative requirement than adding an unconstrained channel multiplier until one secretion percentage is reproduced.

\subsection{Conclusion}

The experimental importance of AE4 survives the reassessment, but the original cation-routing explanation does not remain uniquely determined. A conservation based model reveals how alkalinity supply enables chloride loading, how compensating transport can conceal AE4 loss, and how plausible channel or source constructions can fail physiological or identification tests. Under explicit onset and matched-supply assumptions, a larger chloride reservoir together with continuing AE4 input gives wild type a sustained export advantage through shared secretory pathways. The result provides partial mechanistic support for chloride availability as a basis of AE4 dependent secretion. Its strengths are the substantial conditional effect and exact mass accounting; its unresolved parts are chronic state formation, native compensation and stimulus specificity.
